\section{Training phases and ownership audits}
\label{sup:ownership}

The frozen candidate generator first produces boxes, query features, and base
statistics from the image and complete expression.  The complete-expression
ranker is trained on positives and sealed before the final two phases.
Admission then updates only its category-eligibility owner and a
supervision-only
auxiliary carrier.  Abstention subsequently updates only its sample-level
owner from detached base statistics.  The auxiliary carrier never enters a
deployed score.

\subsection{Exact decision surfaces and objectives}
\label{sup:objectives}

\paragraph{Base score and Ranking.}
Let $h_i$ be final decoder query $i$, let $\mathcal T_e$ be the scored tokens
of the complete expression, and let $\ell_{it}$ be the frozen token logit.  The
base score used by both learned score owners is
\begin{equation}
 s_i=|\mathcal T_e|^{-1}\sum_{t\in\mathcal T_e}\sigma(\ell_{it}).
 \label{eq:supp-base-score}
\end{equation}
Ranking detaches $h_i$ and $s_i$, applies LayerNorm to $h_i$, concatenates
$s_i$, and uses two 128-D GELU layers plus a linear output to predict
$\delta_i^R$; hence $r_i=s_i+\delta_i^R$.  Training uses positive images only.
For labels $y_i=\mathbb 1[\operatorname{IoU}(b_i,b^\star)\geq0.5]$, define
$G(x,y)=\max_{y_i=1}x_i-\max_{y_i=0}x_i$.  A row contributes only if both
labels occur.  If the frozen top-1 is wrong, its loss is
$T_R\operatorname{softplus}((m_R-G(r,y))/T_R)$; otherwise it is
$[[G(s,y)-\epsilon_R]_+-G(r,y)]_+$.  The two row classes are averaged
separately, then summed with
$\lambda_R Q^{-1}\sum_i(\delta_i^R)^2$.  We use
$(T_R,m_R,\epsilon_R,\lambda_R)=(0.1,0.05,0.02,10^{-3})$.

\paragraph{Admission.}
The raw score is
$a_i=\alpha\cos(W_qh_i,P(u))$.  The frozen cue encoder is a patch backbone for
\arrowv, phrase-token mean pooling from the language encoder for \arrowt, and
a fixed dense Rademacher sentinel for \arrown; all feed the same trainable
source projection $P$ and query projection $W_q$, while $\alpha$ is frozen.
For valid mask $\mathcal M$, the deployed standardization is
\begin{equation}
 \begin{aligned}
 \mu&=\frac{1}{n}\sum_{i\in\mathcal M}a_i,\\
 \sigma_a&=\sqrt{\max\!\left(
 \frac{1}{n}\sum_{i\in\mathcal M}(a_i-\mu)^2,10^{-6}\right)},\\
 \tilde a_i&=\operatorname{clip}((a_i-\mu)/\sigma_a,-5,5).
 \end{aligned}
 \label{eq:supp-admission-standardize}
\end{equation}
The auxiliary carrier is a LayerNorm--MLP $3\!\to\!64\!\to\!64\!\to\!1$
with GELUs, applied to
$[\tilde a_i,\tilde r_i,\tilde a_i\tilde r_i]$; it forms
$r_i^{\rm aux}=r_i+\delta_i^A$.  Here $r_i$ is detached and $\tilde r_i$ uses
the same standardization as Eq.~\eqref{eq:supp-admission-standardize}.

Every same-category instance $k$ supplies
$\mathcal P_k=\{i:\operatorname{IoU}(b_i,b_k)\geq0.5\}$, while
$\mathcal N=\{i:\max_k\operatorname{IoU}(b_i,b_k)<0.3\}$.  For instances
with both sets available, the raw-score category margin is
\begin{equation}
 \begin{aligned}
 p_k&=T_A\!\left[\log\!\sum_{i\in\mathcal P_k}e^{a_i/T_A}
       -\log|\mathcal P_k|\right],\\
 \ell_k&=T_A\operatorname{softplus}
       \left(\frac{m_A-p_k+\max_{i\in\mathcal N}a_i}{T_A}\right).
 \end{aligned}
 \label{eq:supp-category-margin}
\end{equation}
The complete Admission objective adds: the repair/preserve loss above on
$r^{\rm aux}$ with positives defined by maximum IoU to any same-category
instance; the preserve term alone for the expression-bound primary instance;
the mean of Eq.~\eqref{eq:supp-category-margin}; and
$10^{-3}Q^{-1}\sum_i(\delta_i^A)^2$.  All three non-regularization weights are
one, and $(T_A,m_A)=(0.1,0.1)$.  For ABI compatibility the current inference
runtime still serializes and computes the carrier, but its output is ignored:
the eligible set is computed solely from $\tilde a$ with margin three, and
eligible queries retain the frozen ranker's order.

\paragraph{Abstention.}
The rejector independently detaches $h_i,s_i$ and maps
$[\operatorname{LN}(h_i),s_i]$ through two 128-D GELU layers to $f_i$.  With
$w_i\propto\exp(s_i/0.01)$ over valid candidates, it concatenates
$\sum_iw_if_i$ with
\begin{equation*}
 \begin{aligned}
 \relax [\max s,\ &\operatorname{mean}(\operatorname{top3}s),
   \operatorname{mean}s,\\
  &\operatorname{std}s,\ s_{(1)}-s_{(2)},
   -\textstyle\sum_iw_i\log w_i/\log|\mathcal M|].
 \end{aligned}
\end{equation*}
A $134\!\to\!128\!\to\!128\!\to\!1$ MLP with GELUs returns one offset $g$
shared by every query.  Thus per-query confidence is $s_i+g$ and the exact
deployed sample score is $c=\max_{i\in\mathcal M}(s_i+g)$.  Neither Admission
scores nor Ranking scores enter this computation.

For paired positive/negative scores $c_j^+,c_j^-$, let $\tau_Q$ be the
detached exact 5th-percentile order statistic of recent positive scores.  The
current distributed batch supplies it until the positive queue reaches 256;
thereafter it uses only the most recent 512 scores from successful updates.
Negative labels come from proposal-covered verification and do not claim that
all frozen queries have been exhaustively annotated.
The main rejection objective is
\begin{align}
 \hat\tau&=\tau_Q+\overline c^+-\sg(\overline c^+),\\
 \mathcal L_-&=\frac{1}{B}\sum_j0.1\,
 \operatorname{softplus}\!\left(\frac{c_j^- -\hat\tau}{0.1}\right),\\
 \mathcal L_+&=\frac{1}{B}\sum_j[\tau_Q-0.02-c_j^+]_+,\\
 \mathcal L_C&=\mathcal L_-+\mathcal L_+ .
 \label{eq:supp-rejection-loss}
\end{align}
The first term keeps every current negative in the loss; negative queue values
do not enter its mean.  The paired-separation term has coefficient zero in the
main model.  The history threshold and queue contents are detached, while the
deployed scores remain live.  Finally, the deployment threshold is selected
on the separate sealed 1,570-row internal calibration set; it is not the
training-queue threshold.

Ownership is checked both structurally and numerically.  Optimizer parameter
identities must equal the phase allowlist; every other tensor is hashed before
and after training; and autograd probes require no path from an Admission loss
to Ranking or Abstention owners, or from an Abstention loss to Admission or
Ranking owners.  Resume receipts additionally bind update counts, optimizer
state, automatic-mixed-precision skips, non-finite events, and frozen hashes.

The critical control separates output count from feature ownership.  One row
uses a shared deployed scalar, a second attaches separate outputs to a shared
trainable feature owner, and a third uses disjoint trainable owners.  An
additional phased schedule is retained as a scheduling diagnostic.  This
original block shows that isolation protects the released frozen-base route;
the capacity-matched and strong-trunk controls below delimit that result rather
than turning it into a universal topology claim.  The evidence does not
establish an independent benefit from phasing once ownership is fixed.

The post-release frozen-base capacity control in
Tab.~\ref{tab:supp-arrow-capacity} reruns both arms with task-specific Adam
states and zero weight decay.  Its Shared-Wide score owner has 83,971
parameters and 83,007 MAC/query, versus 83,969 and 82,944 for Isolated, while
giving either shared duty a 163-d representation rather than 128-d.  Negative
gradient cosines occur in all three Shared-Wide trajectories, yet the paired
Strict FPR95 interval crosses zero.  We therefore retain the result as direct
evidence that shared optimization can damage Ranking, not as evidence of a
capacity-controlled Abstention advantage.

% Generated by paper/scripts/make_main_tables.py.
\begin{table*}[t]
\centering
\caption{ARROW-U2 capacity/optimizer-matched ownership control. Shared-Wide and
Isolated differ by two score-owner parameters; both use two task-specific
AdamW states, zero weight decay, and fixed U150 endpoints.}
\label{tab:supp-arrow-capacity}
\scriptsize
\setlength{\tabcolsep}{4pt}
\begin{tabular}{llrrrrrr}\toprule
Owner & Seed & Test5 & Strict FPR95 & AUROC & AUPR & Grad. cosine & Neg. cosine (\%)\\\midrule
Shared-Wide & 17 & 68.720 & 48.006 & 84.551 & 82.911 & +0.200 & 15.5 \\
Shared-Wide & 42 & 69.818 & 48.154 & 84.421 & 82.816 & +0.145 & 25.0 \\
Shared-Wide & 73 & 68.462 & 47.760 & 84.268 & 82.715 & +0.203 & 18.9 \\
Isolated replay & 17 & 74.290 & 46.972 & 85.059 & 83.419 & zero & 0.0 \\
Isolated replay & 42 & 74.349 & 47.415 & 84.700 & 82.993 & zero & 0.0 \\
Isolated replay & 73 & 74.294 & 47.464 & 84.469 & 82.838 & zero & 0.0 \\
\bottomrule\end{tabular}
\parbox{0.98\textwidth}{\scriptsize The paired Test5 gain of Isolated is
5.311 points (95\% image-cluster CI [4.875, 5.768]).  The paired Strict
FPR95 reduction is 0.689 points (CI [$-0.382$, 1.337], one-sided
$p=0.139$), so the preregistered joint gate does not pass.  Every bootstrap
replicate recomputes each arm/seed's positive q05 threshold.}
\end{table*}


The strong-trunk transfer adds a capacity control that the original block did
not provide.  Shared-128 matches each isolated task tower's width but uses
fewer total parameters.  Shared-Wide instead matches Isolated's total
parameters and MACs while giving either task a wider representation.  All
three learned arms use two task-specific AdamW states and zero weight decay;
the Shared optimizers intentionally maintain independent moments for the same
trunk tensors.  This removes both total capacity and repeated weight decay as
alternative explanations.  The resulting Shared-Wide point estimate is better
than Isolated on Strict-TN2031, so the transfer supports safe route separation
but not universal optimization superiority.

The continuation stress test in
Tab.~\ref{tab:supp-strong-e6-ownership} freezes supplied PosCtrl and TN10 e6
trunks and reruns only Shared-Wide versus Isolated.  It reuses the exact e5
schedule bytes and fixed endpoint.  TN10 increases the negative cosine tail,
but the two routes remain tied on REC and Shared-Wide is significantly better
on Strict-TN2031.  This separates the presence of conflicting probe gradients
from evidence that hard isolation improves deployment.

% Generated by paper/scripts/make_main_tables.py.
\begin{table*}[t]
\centering
\caption{Complete Admission training grid. Val3 is validation-only mechanism evidence; Test5 appears only for the preregistered A5--A1 contrast.}
\label{tab:supp-admission-grid}
\label{tab:supp-training}
\scriptsize
\setlength{\tabcolsep}{3pt}
\begin{tabular}{lcccccc}\toprule
ID & Surface & Aux. & Category-complete & Preserve & Val3 macro$\uparrow$ & Test5$\uparrow$\\\midrule
A0 & Frozen & Frozen & Off & Off & -- & -- \\
A1 & Train & Frozen & On & On & 73.07 $\pm$ 0.21 & 73.54 \\
A2 & Frozen & Train & On & On & -- & -- \\
A3 & Train & Train & On & Off & 73.93 $\pm$ 0.13 & -- \\
A4 & Train & Train & Off & On & 73.80 $\pm$ 0.04 & -- \\
A5 & Train & Train & On & On & 73.89 $\pm$ 0.10 & 74.24 \\
\bottomrule\end{tabular}
\parbox{0.97\textwidth}{\scriptsize Values are three-seed mean $\pm$ sample SD unless only a held-out mean is shown. A2's training-only auxiliary is never deployed and its three deployable checkpoints are bitwise equal to A0. The only sealed static-route score is a seed-17 Val3 micro value of 70.25\%; it is not inserted into the three-seed macro column.}

\end{table*}

\begin{table*}[t]
\centering
\caption{Complete clean Rejection-objective grid and allowed data-provenance controls. Calibration rows are mechanism evidence. Strict-TN2031 was opened only for preregistered C3--C2; no row-specific checkpoint or milestone was selected.}
\label{tab:supp-rejection-grid}
\scriptsize
\setlength{\tabcolsep}{2.3pt}
\begin{tabular}{llcccccc}\toprule
ID & Objective & Queue? & Trust? & Margin? & Cal. FPR95$\downarrow$ & Pair win$\uparrow$ & Strict$\downarrow$\\\midrule
C0 & Identity & No & No & No & -- & -- & -- \\
C1 & Batch negative & No & No & No & 58.81 $\pm$ 2.75 & 75.33 $\pm$ 2.12 & -- \\
C2 & Negative FPR & Yes & No & No & 53.61 $\pm$ 0.37 & 81.30 $\pm$ 0.18 & 46.97 \\
C3 & Negative FPR & Yes & Yes & No & 53.65 $\pm$ 0.91 & 81.42 $\pm$ 0.13 & 47.02 \\
C4 & Negative FPR & Yes & Yes & Yes & 53.86 $\pm$ 0.70 & 81.63 $\pm$ 0.19 & -- \\
\bottomrule\end{tabular}
\vspace{4pt}

\textbf{Data provenance controls}\par\smallskip
\begin{tabular}{llllll}\toprule
ID & Negative source & Verification scope & Cal. FPR95$\downarrow$ & Pair win$\uparrow$ & Role\\\midrule
D0 & None (positive-only) & None & -- & -- & No-TN control \\
D1 & Edited text & Unverified & 81.99 $\pm$ 0.27 & 62.18 $\pm$ 0.18 & Stress control \\
D3 & Proposal-covered pairs & Proposal-covered & 53.65 $\pm$ 0.91 & 81.42 $\pm$ 0.13 & Main clean data \\
\bottomrule\end{tabular}
\parbox{0.97\textwidth}{\scriptsize C0/D0 have no independently trained absolute rejector, hence ``--''. D1 is an unverified weak-data stress control; D3 is the proposal-covered clean source used by C3. Historical rule-swap and approximate matched rows are excluded from manuscript tables.}
\end{table*}

\begin{table*}[t]
\centering
\caption{Complete ownership and scheduling block. Val3/calibration columns are mechanism evidence; held-out cells occur only for the preregistered O2--O0 and O3--O2 contrasts.}
\label{tab:supp-ownership-grid}
\scriptsize
\setlength{\tabcolsep}{2.4pt}
\begin{tabular}{lllccccc}\toprule
ID & Trainable topology & Schedule & Val3$\uparrow$ & Cal. FPR95$\downarrow$ & Pair win$\uparrow$ & Test5$\uparrow$ & Strict$\downarrow$\\\midrule
O0 & Shared scalar & Interleaved & 73.94 $\pm$ 0.04 & 57.32 $\pm$ 0.00 & 80.38 $\pm$ 0.06 & 74.30 & 50.11 \\
O1 & Separate heads/shared feature & Interleaved & 73.87 $\pm$ 0.05 & 57.22 $\pm$ 0.16 & 79.64 $\pm$ 0.10 & -- & -- \\
O2 & Exclusive owners & Interleaved & 73.89 $\pm$ 0.01 & 53.38 $\pm$ 0.32 & 81.17 $\pm$ 0.32 & 74.30 & 47.30 \\
O3 & Exclusive owners & Phased & 73.89 $\pm$ 0.10 & 53.65 $\pm$ 0.91 & 81.42 $\pm$ 0.13 & 74.24 & 47.02 \\
\bottomrule\end{tabular}
\parbox{0.97\textwidth}{\scriptsize O0 has one shared deployed scalar; O1 separates outputs but retains a shared trainable feature; O2/O3 use exclusive owners. O2 establishes the ownership effect. The O3--O2 strict contrast is not significant, so phasing is not claimed as independently beneficial.}
\end{table*}


Admission trajectories use 100 updates, physical batch 56, seeds 17/42/73,
learning rate $3\times10^{-4}$, weight decay $10^{-4}$, gradient clipping at
0.1, and an initial AMP scale of 8192.  Abstention trajectories use 50 updates
and physical batch 8 from the seed-matched frozen Admission parent.  Update
counts are fixed across rows rather than selected separately.  The
current-batch, queue, trust, paired-margin, and unverified-negative rows in
Tab.~\ref{tab:supp-training} are calibration mechanisms; only preregistered
held-out contrasts support paper claims.  In particular, positive trust has no
standalone held-out causal gain, and the unverified-negative row is not treated
as a clean data-scope intervention.
